\chapter{Diseases of the Eyes}
\label{ch:eyes}
The human eye is a complex sensory organ responsible for vision, capable of detecting light across a wide range of intensities and resolving fine detail. Disease of the eye can arise from structural abnormalities, infection, inflammation, degeneration, trauma, or systemic conditions manifesting ocularly. Eye diseases range from common refractive errors to blinding conditions such as glaucoma and age-related macular degeneration.

%-------------------------------------------------------------
\section{Cataract}
%-------------------------------------------------------------

%-------------------------------------------------------------

%-------------------------------------------------------------

%-------------------------------------------------------------

Cataracts are characterized by progressive opacification of the crystalline lens, representing the leading cause of reversible blindness globally. Opacification results from oxidative stress, protein aggregation, metabolic derangements, or mechanical trauma disrupting lens fiber transparency. Surgical phacoemulsification with intraocular lens implantation remains the gold standard treatment with outstanding visual restoration outcomes.

%-------------------------------------------------------------

%-------------------------------------------------------------

%-------------------------------------------------------------

%-------------------------------------------------------------

\label{sec:Cataract}
%-------------------------------------------------------------%-------------------------------------------------------------

\subsection{Age-Related Cataract}
\index{Age-Related Cataract}
\label{sec:Age-Related Cataract}
Age-related cataract is an opacification of the crystalline lens and the leading cause of preventable blindness worldwide, resulting from cumulative oxidative damage to lens proteins. The condition results from aggregation of damaged lens proteins that scatter light. Classification by location includes nuclear (yellowing and hardening of the lens nucleus), cortical (peripheral wedge-shaped opacities), and posterior subcapsular (opacity near the posterior capsule, often associated with steroid use). Clinical features include painless progressive blurring of vision, glare sensitivity, decreased color perception, and dimming of vision. Diagnosis is by slit-lamp examination. Management is surgical: phacoemulsification with implantation of an intraocular lens. Prognosis is excellent, with cataract surgery being one of the most commonly performed and cost-effective surgical procedures globally.

\subsection{Congenital Cataract}
\index{Congenital Cataract}
\label{sec:Congenital Cataract}
Congenital cataract is an opacification of the lens present at or shortly after birth, affecting 1--6 per 10,000 live births. The condition may be idiopathic, hereditary (autosomal dominant most common), or associated with systemic conditions (Down syndrome, congenital rubella, galactosemia, Lowe syndrome, persistent fetal vasculature). The condition results from abnormal lens development or metabolic insult in utero. Clinical features include lens opacity present at birth or in early childhood. Diagnosis is by slit-lamp examination. Management involves urgent surgical removal, ideally within the first few weeks of life for unilateral cases, followed by postoperative visual rehabilitation with patching and corrective lenses. Prognosis depends on timing of intervention and compliance with visual rehabilitation.

%-------------------------------------------------------------
\section{Corneal Diseases}
%-------------------------------------------------------------

%-------------------------------------------------------------

%-------------------------------------------------------------

%-------------------------------------------------------------

Corneal diseases encompass infectious, dystrophic, inflammatory, and degenerative disorders affecting the anterior transparent window of the eye. Pathology compromising corneal avascular clarity or curvature disrupts optical light refraction, causing significant visual impairment and pain. Management ranges from topical antimicrobial and anti-inflammatory therapies to lamellar or penetrating keratoplasty corneal transplantation.

%-------------------------------------------------------------

%-------------------------------------------------------------

%-------------------------------------------------------------

%-------------------------------------------------------------

\label{sec:Corneal Diseases}
%-------------------------------------------------------------%-------------------------------------------------------------

\subsection{Fuchs' Endothelial Dystrophy}
\index{Fuchs' Endothelial Dystrophy}
\label{sec:Fuchs' Endothelial Dystrophy}
Fuchs' endothelial corneal dystrophy is a bilateral, progressive disorder of the corneal endothelium characterized by guttae, endothelial cell loss, and corneal edema; it is autosomal dominant and more common in women. The condition results from progressive endothelial cell dysfunction and loss. Early disease is asymptomatic; as the disease progresses, clinical features include morning blurring that improves through the day, glare, and reduced visual acuity; advanced disease causes persistent corneal edema and bullous keratopathy. Diagnosis is by slit-lamp examination and specular microscopy showing reduced endothelial cell density. Management includes hypertonic saline drops for mild edema and endothelial keratoplasty (DMEK or DSAEK) for advanced disease. Prognosis is excellent with surgical intervention.

\subsection{Keratitis}
\index{Keratitis}
\label{sec:Keratitis}
Keratitis is an inflammation of the cornea that may be infectious or non-infectious. Infectious keratitis can be bacterial, viral (herpes simplex), fungal, or amoebic (\textit{Acanthamoeba}). Bacterial keratitis, the most common type, results from corneal infection, often associated with contact lens use or trauma. Clinical features include pain, redness, tearing, decreased vision, and a corneal infiltrate or ulcer; herpes simplex keratitis may present as epithelial (dendritic ulcer) or stromal disease. Diagnosis involves corneal scraping for Gram stain, culture, and sensitivity. Management requires frequent topical antimicrobial agents appropriate for the suspected organism. Prognosis depends on the causative organism and timeliness of treatment; corneal scarring may require transplantation.

\subsection{Keratoconus}
\index{Keratoconus}
\label{sec:Keratoconus}
Keratoconus is a progressive, non-inflammatory ectatic disorder of the cornea characterized by progressive thinning and steepening, resulting in irregular astigmatism and reduced visual acuity; onset is typically in adolescence or early adulthood. Risk factors include eye rubbing, family history, and associated systemic conditions (Ehlers-Danlos syndrome, Marfan syndrome, Down syndrome). The condition results from structural weakening of the corneal stroma. Clinical features include progressive visual distortion and decreased visual acuity; slit-lamp findings include Munson's sign, Vogt's striae, and corneal iron deposition. Diagnosis is based on corneal topography and tomography. Management progresses from spectacles and rigid gas-permeable contact lenses to corneal crosslinking to halt progression, intrastromal corneal ring segments, and ultimately keratoplasty in advanced cases. Prognosis is variable; corneal crosslinking effectively halts progression.

%-------------------------------------------------------------
\section{Dry Eye Disease}
%-------------------------------------------------------------

%-------------------------------------------------------------

%-------------------------------------------------------------

%-------------------------------------------------------------

%-------------------------------------------------------------

%-------------------------------------------------------------

%-------------------------------------------------------------

%-------------------------------------------------------------

\label{sec:section:Dry Eye Disease}
%-------------------------------------------------------------%-------------------------------------------------------------

\subsection{Dry Eye Disease}
\index{Dry Eye Disease}
\label{sec:Dry Eye Disease}
Dry eye disease is a multifactorial disorder of the ocular surface characterized by loss of homeostasis of the tear film, affecting approximately 5--30\% of the global population, with higher prevalence in the elderly, women, and those in arid or windy environments. The condition results from tear film instability due to aqueous-deficient (reduced lacrimal gland production, as in Sj\"{o}gren's syndrome, aging, or medication side effects) or evaporative causes (meibomian gland dysfunction, lid abnormalities, or environmental factors). Clinical features include dryness, grittiness, burning, foreign body sensation, photophobia, and blurred vision; signs include reduced tear break-up time (less than 10 seconds), Schirmer's test less than 5 mm in 5 minutes, corneal staining, and meibomian gland dropout. Diagnosis is clinical, supported by tear osmolarity testing and inflammatory markers. Management is stepwise: artificial tears, punctal plugs, topical cyclosporine A or lifitegrast, secretagogues, warm compresses and lid hygiene, and environmental modifications. Prognosis is variable, with most patients achieving symptom control with appropriate management.

%-------------------------------------------------------------
\section{Glaucoma}
%-------------------------------------------------------------

%-------------------------------------------------------------

%-------------------------------------------------------------

%-------------------------------------------------------------

Glaucoma is a group of progressive optic neuropathies characterized by characteristic structural changes to the optic nerve head and progressive visual field loss. Elevated intraocular pressure resulting from impaired aqueous humor outflow dynamics represents the primary modifiable risk factor for retinal ganglion cell apoptosis. Early therapeutic lowering of intraocular pressure via topical medications, laser trabeculoplasty, or incisional surgery halts irreversible visual loss.

%-------------------------------------------------------------

%-------------------------------------------------------------

%-------------------------------------------------------------

%-------------------------------------------------------------

\label{sec:Glaucoma}
%-------------------------------------------------------------%-------------------------------------------------------------

\subsection{Glaucoma in Children}
\index{Glaucoma in Children}
\label{sec:Glaucoma in Children}
Primary congenital glaucoma is a condition resulting from developmental abnormality of the trabecular meshwork (goniodysgenesis), causing elevated intraocular pressure from birth or early childhood. The condition results from abnormal development of the aqueous outflow pathways. Clinical features include epiphora, blepharospasm, photophobia, buphthalmos, and corneal edema; juvenile open-angle glaucoma presents in older children and adolescents. Diagnosis is based on examination under anesthesia and tonometry. Management involves early surgical intervention (goniotomy or trabeculotomy); secondary childhood glaucoma may result from aphakia, trauma, uveitis, or steroids. Prognosis depends on early intervention to prevent irreversible optic nerve damage and amblyopia.

\subsection{Primary Angle-Closure Glaucoma}
\index{Primary Angle-Closure Glaucoma}
\label{sec:Primary Angle-Closure Glaucoma}
Primary angle-closure glaucoma is an ophthalmic emergency that occurs when the peripheral iris physically obstructs the trabecular meshwork, blocking aqueous humor drainage and causing a rapid rise in intraocular pressure; it is more common in hyperopic eyes with shallow anterior chambers. The condition results from anatomical predisposition leading to iridotrabecular contact. Clinical features of acute angle closure include severe eye pain, headache, nausea, blurred vision, halos around lights, a mid-dilated fixed pupil, and markedly elevated IOP; chronic angle closure may be asymptomatic until advanced. Diagnosis is based on tonometry and gonioscopy. Management of acute angle closure includes immediate topical and systemic IOP-lowering agents followed by laser peripheral iridotomy; fellow eyes should undergo prophylactic iridotomy. Prognosis is excellent with prompt treatment.

\subsection{Primary Open-Angle Glaucoma}
\index{Primary Open-Angle Glaucoma}
\label{sec:Primary Open-Angle Glaucoma}
Primary open-angle glaucoma is the most common form of glaucoma and a leading cause of irreversible blindness worldwide, characterized by progressive optic nerve damage and visual field loss, typically associated with elevated intraocular pressure. The condition results from increased resistance to aqueous humor outflow through the trabecular meshwork despite an open angle between the iris and cornea. The disease is painless and insidious, often remaining undetected until advanced visual field loss. Clinical features include progressive visual field loss, often beginning peripherally. Diagnosis involves tonometry, gonioscopy, optic disc evaluation (cup-to-disc ratio), automated perimetry, and optical coherence tomography. Management includes topical hypotensive medications (prostaglandin analogs, beta-blockers, alpha-agonists, carbonic anhydrase inhibitors), laser trabeculoplasty, and trabeculectomy or tube shunt surgery for refractory cases. Prognosis is variable, with treatment slowing but not reversing disease progression.

%-------------------------------------------------------------
\section{Inflammatory and Autoimmune Eye Diseases}
%-------------------------------------------------------------

%-------------------------------------------------------------

%-------------------------------------------------------------

%-------------------------------------------------------------

%-------------------------------------------------------------

%-------------------------------------------------------------

%-------------------------------------------------------------

%-------------------------------------------------------------

\label{sec:Inflammatory and Autoimmune Eye Diseases}
%-------------------------------------------------------------%-------------------------------------------------------------

\subsection{Scleritis}
\index{Scleritis}
\label{sec:Scleritis}
Scleritis is a severe inflammation of the sclera that is potentially sight-threatening due to complications such as glaucoma, cataract, and posterior segment inflammation; over 50\% of cases are associated with systemic autoimmune disease, particularly rheumatoid arthritis, granulomatosis with polyangiitis, and relapsing polychondritis. The condition results from immune-mediated inflammation of the sclera. Clinical features include deep, boring eye pain, redness, and decreased vision; anterior scleritis may be diffuse, nodular, or necrotizing; posterior scleritis can mimic orbital tumors. Diagnosis is based on clinical examination and imaging when indicated. Management requires systemic corticosteroids and immunosuppressive agents. Prognosis depends on the severity and underlying systemic disease.

\subsection{Uveitis}
\index{Uveitis}
\label{sec:Uveitis}
Uveitis is an intraocular inflammation involving the uveal tract (iris, ciliary body, choroid), classified by anatomical location as anterior uveitis (iritis or iridocyclitis), intermediate uveitis (pars planitis), posterior uveitis (choroiditis), or panuveitis. The condition results from infectious or autoimmune causes, including HLA-B27-associated conditions (ankylosing spondylitis, reactive arthritis), Beh\c{c}et's disease, sarcoidosis, and infections (herpes, syphilis, tuberculosis). Clinical features of anterior uveitis include eye pain, redness, photophobia, and hypopyon; intermediate uveitis features vitreous cells and floaters; posterior uveitis presents with floaters, visual field defects, or decreased vision. Diagnosis is based on slit-lamp and dilated fundus examination. Management uses topical, periocular, or systemic corticosteroids and steroid-sparing immunosuppressive agents for chronic disease; biologics (anti-TNF agents) are used for refractory cases. Prognosis depends on the underlying cause and severity.

%-------------------------------------------------------------
\section{Ocular Surface Disorders}
%-------------------------------------------------------------

%-------------------------------------------------------------

%-------------------------------------------------------------

%-------------------------------------------------------------

%-------------------------------------------------------------

%-------------------------------------------------------------

%-------------------------------------------------------------

%-------------------------------------------------------------

\label{sec:Ocular Surface Disorders}
%-------------------------------------------------------------%-------------------------------------------------------------

\subsection{Conjunctivitis}
\index{Conjunctivitis}
\label{sec:Conjunctivitis}
Conjunctivitis is an inflammation of the conjunctiva and the most common eye disease worldwide, classified as infectious (viral, bacterial), allergic, or chemical or irritant. Viral conjunctivitis, most commonly caused by adenovirus, results from viral infection of the conjunctiva. Bacterial conjunctivitis results from bacterial infection (S. aureus, S. pneumoniae, H. influenzae, N. gonorrhoeae, C. trachomatis). Clinical features of viral conjunctivitis include acute onset, watery discharge, and preauricular lymphadenopathy; bacterial conjunctivitis presents with purulent discharge and eyelids matted shut; allergic conjunctivitis features bilateral itching, tearing, and stringy mucus. Diagnosis is clinical, with Gram stain and culture used for hyperacute or bacterial cases. Management is etiologic: supportive for viral, topical antibiotics for bacterial, and topical antihistamine or mast cell stabilizers for allergic; gonococcal conjunctivitis requires intramuscular ceftriaxone. Prognosis is excellent with appropriate treatment; untreated gonococcal conjunctivitis can lead to corneal perforation within 24 hours.

%-------------------------------------------------------------
\section{Refractive Errors}
%-------------------------------------------------------------

%-------------------------------------------------------------

%-------------------------------------------------------------

%-------------------------------------------------------------

Refractive errors occur when the optical power of the cornea and crystalline lens fails to focus incoming parallel light rays precisely onto the retina. Mismatches between focal length and axial length produce blurred vision at near or distant ranges. Optical correction via spectacles, contact lenses, or keratorefractive laser procedures restores clear visual acuity.

%-------------------------------------------------------------

%-------------------------------------------------------------

%-------------------------------------------------------------

%-------------------------------------------------------------

\label{sec:Refractive Errors}
%-------------------------------------------------------------%-------------------------------------------------------------

\subsection{Astigmatism}
\index{Astigmatism}
\label{sec:Astigmatism}
Astigmatism is a refractive error caused by irregular curvature of the cornea or lens, resulting in different refractive powers in different meridians. The condition results from irregular corneal or lenticular curvature. Clinical features include blurred or distorted vision at both near and distance. Most astigmatism is corneal and regular. Diagnosis is based on refraction testing. Management involves correction with cylindrical lenses (glasses or contact lenses) or refractive surgery; irregular astigmatism may require rigid gas-permeable contact lenses or scleral lenses. Prognosis is excellent with appropriate correction.

\subsection{Hyperopia}
\index{Hyperopia}
\label{sec:Hyperopia}
Hyperopia is a refractive error that occurs when the axial length of the eye is too short or the refractive power is insufficient, causing light to focus behind the retina. The condition results from insufficient refractive power relative to axial length. Mild hyperopia may be accommodated for, especially in young patients. Clinical features include blurred near vision, eye strain, headaches, and in children, esotropia and amblyopia. Diagnosis is based on refraction testing. Management involves correction with convex (plus) lenses; severe hyperopia in childhood requires early correction to prevent amblyopia and strabismus. Prognosis is excellent with correction.

\subsection{Myopia}
\index{Myopia}
\label{sec:Myopia}
Myopia is the most common refractive error worldwide, affecting approximately 30\% of the global population and approaching 80--90\% in some East Asian countries. The condition results from an axial length of the eye that is too long relative to the refractive power of the cornea and lens, causing light to focus in front of the retina. High myopia (sphere $-6.00$ diopters or greater) increases the risk of retinal detachment, choroidal neovascularization, glaucoma, and cataract. Clinical features include blurred distance vision with preserved near vision. Diagnosis is based on refraction testing. Management involves correction with concave (minus) lenses, contact lenses, or refractive surgery (LASIK, PRK, ICL); atropine eye drops and orthokeratology lenses slow myopia progression in children. Prognosis is excellent with correction.

\subsection{Presbyopia}
\index{Presbyopia}
\label{sec:Presbyopia}
Presbyopia is the age-related loss of accommodation, resulting from hardening of the crystalline lens and decreased ciliary muscle function; it is universal, typically becoming symptomatic after age 40. The condition results from age-related changes in the lens and ciliary muscle. Clinical features include difficulty reading small print, especially in dim light, and the need to hold reading material farther away. Diagnosis is based on clinical history and refraction testing. Management includes reading glasses, bifocal or progressive lenses, monovision contact lens arrangements, and multifocal intraocular lenses; surgical options include corneal inlays, scleral implants, and laser procedures. Prognosis is excellent with appropriate correction.

%-------------------------------------------------------------
\section{Retinal Diseases}
%-------------------------------------------------------------

%-------------------------------------------------------------

%-------------------------------------------------------------

%-------------------------------------------------------------

Retinal diseases involve pathologic processes affecting the neurosensory retina, retinal pigment epithelium, and retinal vascular architecture. Damage to these essential structures disrupts phototransduction pathways, leading to severe vision loss or permanent scotomas. Therapies include intravitreal anti-VEGF microinjections, retinal photocoagulation laser treatment, and vitreoretinal surgical reconstruction.

%-------------------------------------------------------------

%-------------------------------------------------------------

%-------------------------------------------------------------

%-------------------------------------------------------------

\label{sec:Retinal Diseases}
%-------------------------------------------------------------%-------------------------------------------------------------

\subsection{Age-Related Macular Degeneration}
\index{Age-Related Macular Degeneration}
\label{sec:Age-Related Macular Degeneration}
Age-related macular degeneration is the leading cause of irreversible central vision loss in individuals over 50 in developed countries, affecting the macula. The condition results from age-related degenerative changes in the retinal pigment epithelium and photoreceptors; risk factors include age, smoking, family history, and genetic variants (complement factor H). Early AMD is characterized by drusen and pigmentary changes; advanced AMD manifests as either dry (atrophic) AMD with geographic atrophy or wet (neovascular) AMD with choroidal neovascularization. Clinical features include progressive central vision loss and distortion. Diagnosis is based on fundus examination and optical coherence tomography. Management of wet AMD involves anti-VEGF intravitreal injections (ranibizumab, aflibercept, bevacizumab); complement inhibitors are under investigation for dry AMD. Prognosis has improved significantly with anti-VEGF therapy for wet AMD.

\subsection{Diabetic Retinopathy}
\index{Diabetic Retinopathy}
\label{sec:Diabetic Retinopathy}
Diabetic retinopathy is a microvascular complication of diabetes mellitus and a leading cause of blindness in working-age adults. The condition results from chronic hyperglycemia damaging retinal capillaries, causing vascular occlusion, non-perfusion, microaneurysms, hemorrhages, hard exudates, and cotton-wool spots (non-proliferative DR). Progression to proliferative DR involves neovascularization driven by retinal ischemia and VEGF upregulation, with risks of vitreous hemorrhage and tractional retinal detachment; diabetic macular edema can occur at any stage. Clinical features include vision loss, which may be asymptomatic until advanced. Diagnosis involves dilated fundus examination and fundus photography. Management includes optimal glycemic control, anti-VEGF intravitreal injections for diabetic macular edema, focal or grid laser photocoagulation, and panretinal photocoagulation for proliferative disease; regular screening is recommended for all diabetic patients. Prognosis depends on glycemic control and stage at diagnosis.

\subsection{Retinal Detachment}
\index{Retinal Detachment}
\label{sec:Retinal Detachment}
Retinal detachment is a separation of the neurosensory retina from the underlying retinal pigment epithelium, constituting an ophthalmic emergency; rhegmatogenous retinal detachment, the most common type, results from a retinal break allowing vitreous fluid to enter the subretinal space. Risk factors include high myopia, prior cataract surgery, trauma, and lattice degeneration. Tractional retinal detachment occurs in proliferative diabetic retinopathy; exudative retinal detachment results from fluid accumulation without a break. Clinical features include sudden onset of floaters, flashing lights, and a curtain-like visual field defect. Diagnosis is based on dilated fundus examination. Management of rhegmatogenous retinal detachment requires surgical repair with pneumatic retinopexy, scleral buckling, or vitrectomy. Prognosis depends on the type, location, and timeliness of repair.

\subsection{Retinitis Pigmentosa}
\index{Retinitis Pigmentosa}
\label{sec:Retinitis Pigmentosa}
Retinitis pigmentosa is a group of inherited retinal dystrophies characterized by progressive degeneration of photoreceptors, primarily rod cells; most cases are inherited (autosomal dominant, autosomal recessive, or X-linked), with over 80 genes identified. The condition results from genetic mutations causing photoreceptor degeneration. Clinical features include night blindness and progressive peripheral visual field constriction (tunnel vision), eventually affecting central vision; fundoscopy reveals bone-spicule pigmentation, waxy pallor of the optic disc, and attenuated retinal vessels. Diagnosis is based on fundus examination and electroretinography showing diminished or absent responses. Management is supportive; vitamin A supplementation may slow progression in some forms; retinal prostheses and gene therapy (voretigene neparvovec for RPE65-mutated RP) offer emerging therapeutic avenues. Prognosis is progressive with no cure currently available.

\subsection{Retinopathy of Prematurity}
\index{Retinopathy of Prematurity}
\label{sec:Retinopathy of Prematurity}
Retinopathy of prematurity is a vasoproliferative retinal disease affecting premature infants, particularly those born before 32 weeks gestation with low birth weight. The condition results from abnormal vascularization of the developing retina, which can lead to tractional retinal detachment and blindness. The disease is classified by zone, stage, and extent of vascular involvement, with plus disease indicating vascular activity. Clinical features depend on disease severity; most mild ROP regresses spontaneously. Diagnosis is based on screening examinations in at-risk premature infants. Management of severe ROP involves anti-VEGF intravitreal injections or laser photocoagulation of the avascular retina. Prognosis is good with timely screening and treatment; long-term follow-up is necessary due to risks of late sequelae including myopia, strabismus, and retinal detachment.

%-------------------------------------------------------------
\section{Strabismus and Amblyopia}
%-------------------------------------------------------------

%-------------------------------------------------------------

%-------------------------------------------------------------

%-------------------------------------------------------------

%-------------------------------------------------------------

%-------------------------------------------------------------

%-------------------------------------------------------------

%-------------------------------------------------------------

\label{sec:Strabismus and Amblyopia}
%-------------------------------------------------------------%-------------------------------------------------------------

\subsection{Amblyopia}
\index{Amblyopia}
\label{sec:Amblyopia}
Amblyopia is reduced best-corrected visual acuity in one or both eyes resulting from abnormal visual development during the critical period of visual maturation (birth to approximately 8 years); it is the most common cause of monocular vision loss in children. The condition results from strabismus, anisometropia, form deprivation (congenital cataract, ptosis), or bilateral high refractive error, leading to failure of normal neural pathway development in the amblyopic eye. Clinical features include reduced visual acuity that cannot be corrected with lenses. Diagnosis is based on visual acuity testing. Management includes corrective lenses, patching or penalization (atropine drops) of the dominant eye, and treatment of the underlying cause. Prognosis is excellent with early intervention.

\subsection{Strabismus}
\index{Strabismus}
\label{sec:Strabismus}
Strabismus is a misalignment of the visual axes of the two eyes that may be horizontal (esotropia, exotropia) or vertical (hypertropia, hypotropia), and may be constant or intermittent, comitant or non-comitant. The condition results from neuromuscular imbalance of the extraocular muscles; onset in childhood may be accommodative, paralytic, or idiopathic; in adults, acquired strabismus may result from cranial nerve palsies (III, IV, VI), thyroid eye disease, myasthenia gravis, or orbital pathology. Clinical features include visible eye misalignment, which can cause amblyopia in children and diplopia in adults. Diagnosis is based on cover testing and ocular motility examination. Management includes corrective lenses, prisms, occlusion therapy, botulinum toxin injection, and surgical realignment of the extraocular muscles. Prognosis is generally good with appropriate treatment.

